\documentclass[11pt]{article}
\usepackage{amsmath,amssymb,amsthm}
\usepackage[margin=1in]{geometry}
\usepackage{hyperref}

\newtheorem{definition}{Definition}[section]
\newtheorem{theorem}[definition]{Theorem}
\newtheorem{lemma}[definition]{Lemma}
\newtheorem{corollary}[definition]{Corollary}
\newtheorem{proposition}[definition]{Proposition}
\newtheorem{remark}[definition]{Remark}

\title{The Baseless Logarithm: A Formal Development}
\author{Formalization of ideas from ``Everything Is Logarithms'' by Alex Kritchevsky}
\date{\today}

\begin{document}
\maketitle

\begin{abstract}
We formalize the notion of a \emph{baseless logarithm} as an abstract algebraic object, independent of any choice of base. We show that baseless logarithms form a torsor under the multiplicative group $\mathbb{R}^{>0}$, and that the standard logarithm $\log_b x$ arises as a projection of this object onto a chosen coordinate system (base). We establish the change-of-base formula as a consequence of this structure, and connect the framework to vectors in differential geometry.
\end{abstract}

\tableofcontents
\newpage

%======================================================================
\section{Motivation}
\label{sec:motivation}
%======================================================================

The logarithm $\log_b x$ is conventionally defined as the solution $y$ to $b^y = x$. This definition ties the logarithm to a specific base $b$, making it difficult to see that the \emph{operation} of taking a logarithm is independent of the choice of base.

We propose a more natural object: the \textbf{baseless logarithm}, written $\log x$, which captures the multiplicative structure of $x$ without reference to any unit. Just as a geometric vector $\mathbf{v}$ exists independent of coordinates, $\log x$ exists independent of base.

%======================================================================
\section{Definitions}
\label{sec:definitions}
%======================================================================

\begin{definition}[Baseless Logarithm]
\label{def:baseless}
The \textbf{baseless logarithm} of a positive real number $x \in \mathbb{R}^{>0}$ is the formal expression
\[
\log x
\]
regarded as an element of an abstract set $\mathcal{L}$ (the \textbf{logarithmic space}).
\end{definition}

\begin{remark}
The expression $\log x$ is not assigned a numerical value. It is an algebraic object, analogous to a geometric vector $\mathbf{v}$ before projection onto coordinates.
\end{remark}

\begin{definition}[Based Logarithm as Projection]
\label{def:based}
Given a base $b \in \mathbb{R}^{>0}$ with $b \neq 1$, the \textbf{based logarithm} is the projection
\[
\log_b x \;=\; \frac{\log x}{\log b} \;\in\; \mathbb{R}.
\]
Here $\log b$ plays the role of a \textbf{unit vector} or \textbf{measuring stick} that converts the abstract object $\log x$ into a real number.
\end{definition}

\begin{example}
\begin{itemize}
\item The ``bits'' unit: $\log_2 x = \frac{\log x}{\log 2}$. Here $\log 2$ is the unit ``1 bit.''
\item The ``nats'' unit: $\ln x = \frac{\log x}{\log e}$. Here $\log e$ is the unit ``1 nat.''
\end{itemize}
\end{example}

%======================================================================
\section{Algebraic Structure of $\mathcal{L}$}
\label{sec:structure}
%======================================================================

We now establish the algebraic properties of the baseless logarithm.

\begin{definition}[Logarithmic Space]
\label{def:logspace}
Let $\mathcal{L}$ be the set of all baseless logarithms. We equip $\mathcal{L}$ with a scalar multiplication by $\mathbb{R}^{>0}$:
\[
\cdot \;:\; \mathbb{R}^{>0} \times \mathcal{L} \to \mathcal{L}, \qquad c \cdot (\log x) \;=\; \log(x^c).
\]
\end{definition}

\begin{proposition}[Multiplicative-to-Additive Isomorphism]
\label{prop:isomorphism}
The map
\[
\log \;:\; (\mathbb{R}^{>0}, \times) \;\to\; (\mathcal{L}, +)
\]
defined by $\log(xy) = \log x + \log y$ is a group homomorphism, where addition in $\mathcal{L}$ is defined by
\[
(\log x) + (\log y) \;=\; \log(xy).
\]
\end{proposition}

\begin{proof}
This is immediate from the definition of logarithms: for any bases $b_1, b_2$,
\[
\frac{\log(xy)}{\log b} \;=\; \frac{\log x + \log y}{\log b} \;=\; \frac{\log x}{\log b} + \frac{\log y}{\log b} \;=\; \log_b x + \log_b y.
\]
Since this holds for all bases $b$, the identity $\log(xy) = \log x + \log y$ holds in $\mathcal{L}$.
\end{proof}

\begin{theorem}[Change of Base as Unit Conversion]
\label{thm:changebase}
For any $b_1, b_2, x \in \mathbb{R}^{>0}$ with $b_1, b_2 \neq 1$:
\[
\log_{b_1} x \;=\; \log_{b_1}(b_2) \cdot \log_{b_2} x.
\]
That is, changing the base is equivalent to multiplying by the \emph{conversion factor} $\log_{b_1}(b_2)$ between the two units.
\end{theorem}

\begin{proof}
\begin{align*}
\log_{b_1}(b_2) \cdot \log_{b_2} x
&= \frac{\log b_2}{\log b_1} \cdot \frac{\log x}{\log b_2} \\
&= \frac{\log x}{\log b_1} \\
&= \log_{b_1} x. \qedhere
\end{align*}
\end{proof}

\begin{remark}
This is precisely analogous to unit conversion in physics:
\[
2 \text{ km} = \frac{2000 \text{ m}}{1000 \text{ m/km}}.
\]
The ``units'' $\log b_1$ and $\log b_2$ play the same role as ``meters'' and ``kilometers.''
\end{remark}

%======================================================================
\section{Torsor Structure}
\label{sec:torsor}
%======================================================================

We now show that $\mathcal{L}$ is not a vector space, but a \textbf{torsor} (principal homogeneous space) under $\mathbb{R}^{>0}$.

\begin{definition}[Torsor]
A \textbf{torsor} under a group $G$ is a set $T$ with a free and transitive action of $G$. That is, for any $t_1, t_2 \in T$, there exists a unique $g \in G$ such that $t_2 = g \cdot t_1$.
\end{definition}

\begin{theorem}[$\mathcal{L}$ is a Torsor under $\mathbb{R}^{>0}$]
\label{thm:torsor}
The logarithmic space $\mathcal{L}$ is a torsor under the multiplicative group $(\mathbb{R}^{>0}, \times)$, with action
\[
c \cdot (\log x) \;=\; \log(x^c), \qquad c \in \mathbb{R}^{>0}.
\]
\end{theorem}

\begin{proof}
We verify the two torsor axioms:

\textbf{Free:} If $c \cdot (\log x) = \log x$ for all $x$, then $\log(x^c) = \log x$, so $x^c = x$ for all $x > 0$, hence $c = 1$.

\textbf{Transitive:} Given $\log x$ and $\log y$ in $\mathcal{L}$, we need $c$ such that $\log(x^c) = \log y$. This gives $x^c = y$, so $c = \frac{\log y}{\log x}$ (which is well-defined for $x \neq 1$; if $x = 1$ then $\log x = 0$ and the action is trivial).
\end{proof}

\begin{remark}
This is analogous to the torsor of geometric vectors: a vector $\mathbf{v}$ exists independent of coordinates, and choosing a basis is like choosing a ``unit'' for measurement.
\end{remark}

%======================================================================
\section{Connection to Vectors}
\label{sec:vectors}
%======================================================================

We now draw the formal correspondence between baseless logarithms and geometric vectors.

\begin{table}[h]
\centering
\begin{tabular}{cc}
\hline
\textbf{Logarithms} & \textbf{Vectors} \\
\hline
$\log x$ (baseless) & $\mathbf{v}$ (geometric) \\
$\log b$ (unit) & $\mathbf{e}$ (basis vector) \\
$\log_b x = \frac{\log x}{\log b}$ & $v_e = \frac{\mathbf{v}}{\mathbf{e}}$ (projection) \\
Change of base & Change of coordinates \\
$b_1 \leftrightarrow b_2$ & $\mathbf{e}_1 \leftrightarrow \mathbf{e}_2$ \\
$\log_{b_1} x = \log_{b_1}(b_2) \cdot \log_{b_2} x$ & $v_{e_1} = \frac{e_2}{e_1} \cdot v_{e_2}$ \\
\hline
\end{tabular}
\caption{Correspondence between logarithmic and vector operations.}
\label{tab:correspondence}
\end{table}

\begin{proposition}[Vector-Logarithm Analogy]
\label{prop:vector-analogy}
The following diagram commutes:
\[
\begin{array}{ccc}
\mathcal{L} & \xrightarrow{\text{choose unit } \log b} & \mathbb{R} \\
\downarrow_{\text{addition}} & & \downarrow_{+} \\
\mathcal{L} & \xrightarrow{\text{choose unit } \log b} & \mathbb{R}
\end{array}
\]
That is, adding baseless logarithms and then projecting to a base is the same as projecting to a base and then adding real numbers.
\end{proposition}

\begin{proof}
\[
\frac{\log(xy)}{\log b} = \frac{\log x + \log y}{\log b} = \frac{\log x}{\log b} + \frac{\log y}{\log b}. \qedhere
\]
\end{proof}

%======================================================================
\section{Partial Logarithms and Projections}
\label{sec:partial}
%======================================================================

The article observes that there is no direct analogue of \emph{partial derivatives} for logarithms---yet various fields have independently invented such operations.

\begin{definition}[$p$-adic Valuation as Logarithmic Projection]
\label{def:padic}
For a prime $p$ and positive integer $n$ with factorization $n = \prod_i p_i^{a_i}$, define
\[
\nu_p(n) \;=\; a_p
\]
as the \textbf{projection} of $\log n$ onto the $\log p$ component:
\[
\log n \;=\; \sum_i a_i \log p_i, \qquad \nu_p(n) = a_p.
\]
\end{definition}

\begin{proposition}[$\nu_p$ is Logarithmic]
\label{prop:nulog}
The $p$-adic valuation satisfies:
\begin{enumerate}
\item $\nu_p(mn) = \nu_p(m) + \nu_p(n)$ (additive in the multiplicative argument).
\item $\nu_p(m/n) = \nu_p(m) - \nu_p(n)$.
\item $\nu_p(p^k) = k$.
\end{enumerate}
\end{proposition}

\begin{proof}
These follow directly from the uniqueness of prime factorization. If $m = p^a \cdot m'$ and $n = p^b \cdot n'$ with $p \nmid m', n'$, then $mn = p^{a+b} \cdot m'n'$ and $m/n = p^{a-b} \cdot m'/n'$.
\end{proof}

\begin{remark}
We may write this schematically as
\[
\nu_p(n) \;=\; \frac{\partial}{\partial (\log p)} \log n
\]
where the ``derivative'' is the projection onto the $\log p$ direction. This mirrors $\frac{\partial f}{\partial x} = \frac{df}{dx}\big|_{\text{$x$-component}}$ in vector calculus.
\end{remark}

%======================================================================
\section{Summary and Open Questions}
\label{sec:summary}
%======================================================================

We have formalized:
\begin{enumerate}
\item The baseless logarithm $\log x$ as an abstract object in $\mathcal{L}$.
\item The based logarithm $\log_b x$ as a projection $\frac{\log x}{\log b}$.
\item Change of base as unit conversion: $\log_{b_1} x = \log_{b_1}(b_2) \cdot \log_{b_2} x$.
\item $\mathcal{L}$ as a torsor under $(\mathbb{R}^{>0}, \times)$.
\item The correspondence with geometric vectors.
\item $p$-adic valuation as a ``partial logarithm.''
\end{enumerate}

\textbf{Open questions:}
\begin{itemize}
\item Can we define a general \emph{partial logarithm} operator for arbitrary multiplicative decompositions?
\item How does this framework extend to complex logarithms and the multi-valued structure?
\item What is the precise relationship between $\mathcal{L}$ and the Lie algebra of $(\mathbb{R}^{>0}, \times)$?
\item Can fractional dimensions (as in \S\ref{sec:motivation}) be given a rigorous algebraic interpretation?
\end{itemize}

\end{document}
